\documentclass[main.tex]{subfiles}

\begin{document}\setWPP
\lsection{Funkce}
\begin{multicols}{2}

\qaw{Vlastnosti a definiční obor funkce}{
Předpis funkce:
$$f(x) = \frac{\cos x}{1 - x^2}$$
Definiční obor:
$$D_f = (-\infty, -1) \cup (-1, 1) \cup (1, \infty)$$
Parita funkce:
$$f(-x) = \frac{\cos(-x)}{1 - (-x)^2} = \frac{\cos x}{1 - x^2} = f(x)$$
Funkce je \texthl{sudá}.
}
\columnbreak
\qaw{Hledání inverzní funkce}{
Předpis funkce:
$$f\colon y = \sqrt{e^{-(1+x)}-1}$$
Inverzní funkce:
$$
\begin{array}{l @{\qquad} r @{} c @{} l}
    f^{-1}\colon & y^2 &=& e^{-(1+x)}-1 \\
    f^{-1}\colon & y^2+1 &=& e^{-(1+x)} \\
    f^{-1}\colon & \ln(y^2+1) &=& -1-x \\
    f^{-1}\colon & -1-\ln(y^2+1) &=& x \\
    f^{-1}\colon & y &=& -1-\ln(x^2+1)
\end{array}
$$
Definiční obory:
$$D_f = (-\infty, -1\rangle$$
$$D_{f^{-1}} = \langle0, \infty)$$
}
\end{multicols}
\qaw{Vysvětli pojmy:}{%
\setlength{\multicolsep}{0pt}% Remove vertical space before and after
\setlength{\columnsep}{2em}% Remove space between columns
\setlength{\leftskip}{0em}%
\SetMCRule{
  custom-line={\path[WPP] (TOP) to [ornament=88] (BOT);}, % Use ornament 88
  width=0.4pt, % Ensure rule is drawn (width > 0pt)
  extend-top=-2pt, % Shrink rule at top by 24pt
  extend-bot=-8pt % Shrink rule at bottom by 8pt
}
\hspace*{2em}
\parbox{\dimexpr\textwidth-2em}{
\begin{multicols}{2}
\textbl{daňová záloha}
-- jedná se o část daně, která se platí ještě \texthl{před ukončením zdaňovacího období}\\
\textbl{plátce daně}
-- ten, \texthl{kdo daň odvádí státu} (u nepřímých daň vybírá, u přímých je zároveň poplatníkem)\\
\textbl{poplatník}
-- ten, \texthl{z jehož zdrojů} se daň \texthl{reálně odvádí} (u DPH zákazník, u daně z příjmů zaměstnanec)\\
\textbl{přeplatek na dani}
-- když státu \texthl{odvedeme víc}, než máme povinnost (například u moc vysokých záloh) a stát nám \texthl{přeplatek vratí}\\
\textbl{daňové přiznání}
-- oficiální dokument, kde je plátce \texthl{povinen} uvést \texthl{všechny skutečnosti} potřebných \texthl{k výpočtu daně}, lživé přiznání je \texthl{trestné}\\
\textbl{správce daně}
-- finanční úřad, který \texthl{vyměřuje výši} daně\\
\columnbreak\\
{\color{antiDP}$\faGreaterThan$}\textbl{předmět daně}
-- všechny příjmy \texthl{podléhající zdanění}, stejně tak \texthl{majetek} či poskytnuté \texthl{služby}\\
{\color{antiDP}$\faGreaterThan$}\textbl{osvobození z daně z příjmů fyzických osob}
-- úlevy stanovené zákonem, které \texthl{vyjímají plátce z povinnosti} daň platit\\
{\color{antiDP}$\faGreaterThan$}\textbl{zdaňovací období}
-- období, za které daň odvádíte, obvykle \texthl{jeden rok} nebo \texthl{jeden měsíc} (u DPH je možno nově i~\texthl{čtvrťletně})
\end{multicols}
}}
\vspace{2em}
\wbreak
\qaw{Jaké možnosti podání daňového přiznání má OSVČ}{%
OSVČ musí podávat přiznání přes \texthl{Datovou Schránku}, má k tomu tři možnosti.\\
\begin{itemize}[label={\color{WPP}\faSlackHash}]
\item \textbl{paušální daň} -- pro neplátce DPH, jedná se o \texthl{kombinaci ZP, SP a daně z příjmů}, vhodná pro menší podnikatele
\item \textbl{výdajový paušál} -- pro základ daně použije \texthl{část svých příjmů} (20 \% -- 80 \% dle sektoru), nemusí tedy dokládat reálné výdaje, což zjednodušuje přiznání
\item \textbl{sledování příjmů a výdajů} -- OSVČ, který vede daňovou evidenci si spočítá daňový základ jako \texthl{rozdíl mezi výdaji a příjmy}, musí tedy řádně vést administrativu, ale platí jen to, co reálně musí
\end{itemize}
}
\wpage
\end{document}